% ============================================
% Introduction
% ============================================
\section{Introduction}
\label{sec:introduction}

% ============================================
% NEW INTRODUCTION (in red) - Do not remove old version below
% ============================================

\textcolor{red}{
% Paragraph 1: The fundamental question - communities as cohesive units
A long-standing question in ecology is whether communities behave as cohesive units or merely as collections of independently acting species \citep{Shapiro1998, West2006}. If communities function as integrated wholes, species within them should experience correlated fates during ecological perturbations, and community-level properties should predict ecological outcomes better than individual species traits. Conversely, if species act independently, community responses would simply reflect the sum of individual species dynamics. This distinction has profound implications for understanding ecosystem stability, predicting responses to environmental change, and designing interventions such as probiotic therapies or ecosystem restoration.
}

\textcolor{red}{
% Paragraph 2: Contrasting approaches and views
Multiple theoretical and empirical approaches have been developed to address this question, often yielding contrasting perspectives. Consumer--resource theory suggests that communities can behave as quasi-units when species share limiting resources, with community-level metrics predicting survival better than individual fitness \citep{Tikhonov2016}. [CITE XX - additional theory]. Empirically, microbial experiments using top-down and bottom-up approaches have revealed coselection mechanisms: dominant taxa can recruit subdominants through cross-feeding, and vice versa, suggesting coordinated community behavior \citep{DiazColunga2022}. Studies of fecal microbiota transplantation (FMT) show that donor-strain engraftment depends jointly on donor and recipient composition, with correlated engraftment outcomes among donor strains consistent with shared selection pressures \citep{Smillie2018}. [CITE YYY - contrasting view]. Despite these advances, a unifying framework that reconciles these perspectives and identifies when communities behave cohesively remains elusive.
}

\textcolor{red}{
% Paragraph 3: Coalescence as a natural test of community cohesion
Community coalescence---the merging of entire communities rather than single-species invasions---provides a natural experimental paradigm to test community cohesion \citep{Rillig2015, Lechon2021}. Coalescence occurs ubiquitously across scales: mixing soil or water samples combines microbial assemblages at the microscale \citep{Liu2024}, river confluences merge freshwater and marine communities \citep{Rocca2020}, and host-associated microbiomes coalesce through fecal transplantation or social contact \citep{Xiao2020}. Critically, coalescence outcomes can reveal whether communities act as cohesive units: if species within a community experience correlated survival or extinction during coalescence, this suggests they function as an integrated whole subject to community-level selection. Alternatively, if species fates are uncorrelated and determined by individual competitive abilities, coalescence outcomes would reflect species-level rather than community-level dynamics. Thus, by analyzing whether coalescence produces community-level selection (one community overwhelming the other) versus independent species sorting, we can directly probe the conditions under which communities behave as cohesive units.
}

\textcolor{red}{
% Paragraph 4: Our work
Here, we combine theoretical modeling with synthetic-microcosm experiments to investigate when communities behave as cohesive units during coalescence and how interspecies interaction strength governs this behavior. We map post-coalescence communities into a similarity space relative to their parental communities and classify three outcome types: \textit{community-level selection (CLS)} where one parental community overwhelms the other, \textit{Mixture} where both parents persist comparably, and \textit{Restructuring} where a novel state emerges distinct from both parents. Using a generalized Lotka--Volterra model, we show that interaction strength determines the frequency of each outcome: weak interactions favor Mixture, while strong interactions promote CLS. We experimentally validate these predictions by tuning nutrient concentration as a control knob for interaction strength. Furthermore, we identify two regimes of community competitiveness underlying CLS: a species-driven regime under strong interactions where dominant taxa determine the winner, and an emergent regime under intermediate interactions where many-body dynamics collectively shape the outcome. Together, our results demonstrate that community-level selection, and its predictability, emerges conditionally from interspecies interactions, providing a quantitative framework linking interaction strength to community cohesion.
}

% ============================================
% ORIGINAL INTRODUCTION (keep for reference)
% ============================================

\textcolor{gray}{[ORIGINAL INTRODUCTION BELOW - TO BE REMOVED LATER]}

In ecology, the invasion of microbial species into an existing community can profoundly alter biodiversity, ecosystem functioning, and habitat resilience \citep{Liu2024, Ratzke2020, Amor2020} and has been widely studied. However, in nature, invasion of entire communities rather than single species to another community is frequent---termed \textit{community coalescence} \citep{Lechon2021, Rillig2015}. Community coalescence occurs in diverse contexts and scale. At the microscale, combining soil or water samples mixes their microbial assemblages \citep{Liu2024}. Aquatic environments provide larger-scale examples: at river confluences, freshwater and marine communities collide and mix \citep{Rocca2020}. Host-associated microbiomes can also coalesce in contexts such as fecal microbiota transplantation (FMT) or daily social contact among animals \citep{Xiao2020}.

Despite its prevalence and importance in community assembly, community coalescence remains poorly understood \citep{Rillig2015, Castledine2020, Rillig2017}. Intracommunity preselection implies that each parental community arrives as an environmentally filtered, adapted ensemble of coexisting species with structured abundance profile. Intercommunity encounters introduce cross-feeding, inhibition, and phage transfer that arise only upon mixing, promoting extinction events or complete restructuring of the interaction network. The interplay between intracommunity preselection and intercommunity encounter generates non-additive assembly in which the realized fitness of taxa is often context dependent at the community scale \citep{Smillie2018, Huet2023}. Consequently, community coalescence is not adequately captured by single-species invasion frameworks, and instead represents a distinct ecological paradigm.

Despite this complexity, a long-standing idea treats communities as cohesive units, with species interacting in structured ways and experiencing correlated selection \citep{Shapiro1998, West2006, Tikhonov2016}. Under this view, coalescence between cohesive communities can realize \textit{community-level selection} where outcomes are determined by the two communities' collective properties and describable in lower-dimensional, tractable terms. Consumer--resource theory emphasizes that communities can behave as quasi-units and that community-level metrics better predict survival than individual metrics \citep{Tikhonov2016}. Empirically, coalescence experiments reveal coselection mediated by cross-feeding mechanism in both top-down and bottom-up manners, where dominant taxa recruit subdominants and vice versa \citep{DiazColunga2022}. In host-associated systems, strain-level analyses after fecal microbiota transplantation (FMT) show that donor-strain engraftment depends jointly on donor and recipient composition, with correlated engraftment outcomes among donor strains consistent with shared selection within the donor community \citep{Smillie2018}. Together, these lines of evidence suggest that communities can act as integrated units during coalescence, yet a general framework to study community cohesion and resulting coalescence dynamics across systems remains lacking.

Here, we study when communities behave as cohesive units during coalescence and how interspecies interactions govern the outcomes in the lens of interaction strength, the very generic metric. Combining theoretical modeling with synthetic-microcosm experiments, we mapped post-coalescence communities into a similarity space relative to their parental communities and classified three outcome types: \textit{CLS} (one parental community overwhelms the other), \textit{Mixture} (both parents persist at comparable levels), and \textit{Restructuring} (a novel state distinct from both parents). In an initial experiment with Base medium, where interspecies interactions are moderate, we observe community-level selection (CLS), as evidenced by frequent dominance of one parental community. A generalized Lotka--Volterra model reproduces this prevalence, indicating that interaction strength can capture the frequency of community-level selection. The model further predicts how interaction strength shapes coalescence outcomes: under weak interactions, Mixture is favored, whereas CLS becomes increasingly common as interactions strengthen. We then experimentally recapitulate this pattern by tuning the nutrient concentration of the medium, a control knob for interspecies interaction strength, and show that decreasing interaction intensity shifts outcomes away from CLS and Restructuring to Mixture, consistent with the predictions. Community-level selection occurs in both Base (intermediate interaction strength) and Nutr+ (strong interaction strength), yet predictability analyses reveal two regimes of community competitiveness: a species-driven regime under strong interactions, where a few dominant taxa largely determine the winner, and an emergent regime under intermediate interactions, where many-body dynamics collectively shape the outcome. Together, these results show that community-level selection, and the degree to which it is predictable, emerges conditionally from interspecies interactions.
